\documentclass[sigconf]{acmart}
\usepackage{amsmath}
\usepackage{booktabs}
\usepackage{hyperref}
\title{NeuroDigital Interfaces and Adaptive Cognitive Networking: A Nanotechnology-Ready Framework for Human-AI Cognitive Interfaces}
\author{Fernando May Fuentes}
\affiliation{\institution{Instituto Polit\'{e}cnico Nacional, UPIITA}\city{Mexico City}\country{Mexico}}
\affiliation{\institution{Beihang University, School of Information and Communication Engineering}\city{Beijing}\country{China}}
\email{fmayf1500@alumno.ipn.mx}
\thanks{ORCID: \href{https://orcid.org/0009-0002-3953-5224}{0009-0002-3953-5224}.}

\begin{document}
\begin{abstract}
NeuroDigital Interfaces (NDI) provide a systems pathway for translating non-invasive neural observations into machine-interpretable cognitive states while preserving uncertainty, privacy, and human agency. This paper presents NDAN, a NeuroDigital Adaptive Network that combines BCI signal preprocessing, probabilistic cognitive-state estimation, adaptive networking, and a governance layer for human-AI interaction. The proposed framework is nanotechnology-ready rather than nanotechnology-validated: nanosensors, flexible electrodes, and low-power edge devices are specified as future sensing substrates, while the current validation layer is software-based. We formalize the observation-to-intent pipeline, define uncertainty-aware state communication, and propose an evaluation protocol covering signal quality, latency, calibration, energy, privacy, and human-in-the-loop utility. The paper avoids claims of direct telepathy or clinical efficacy and instead establishes a reproducible architecture and translational roadmap for safe cognitive interfaces.
\end{abstract}
\keywords{neurodigital interfaces, BCI, cognitive networking, nanosensors, human-AI interaction, neuro-rights}
\maketitle

\section{Introduction}
Human-AI interaction is increasingly constrained by the bandwidth and ambiguity of conventional input channels. NeuroDigital Interfaces can complement these channels by estimating task-relevant cognitive states from non-invasive measurements such as EEG or fNIRS. However, the engineering challenge is not simply signal classification: a useful interface must quantify uncertainty, adapt to changing users and environments, communicate safely with distributed agents, and preserve autonomy.

This paper introduces NDAN as a system architecture rather than a claim of mind reading. The nanotechnology connection is defined as a future sensing and packaging layer involving nanosensors, flexible electrode interfaces, and low-power edge electronics. Those components require independent hardware and biomedical validation and are not presented as experimentally demonstrated here.

\section{System Architecture}
NDAN is organized into five layers:
\begin{enumerate}
\item \textbf{Sensing:} neural, physiological, and contextual observations;
\item \textbf{Signal processing:} filtering, artifact detection, feature extraction, and calibration;
\item \textbf{Cognitive state estimation:} probabilistic estimates of attention, workload, fatigue, and task intent;
\item \textbf{Adaptive networking:} state-aware communication among human, edge, and AI agents;
\item \textbf{Governance:} consent, provenance, access control, uncertainty disclosure, and reversibility.
\end{enumerate}

The resulting loop is:
\begin{equation}
 y_t \rightarrow \phi(y_t) \rightarrow p(z_t|\phi) \rightarrow m_t \rightarrow a_t \rightarrow y_{t+1},
\end{equation}
where $y_t$ is a sensor observation, $z_t$ is a latent cognitive state, $m_t$ is a privacy-preserving message, and $a_t$ is an adaptive system action.

\section{Nanotechnology-Ready Interface}
The architecture identifies three interfaces where nanoscale engineering may provide future benefits:
\begin{itemize}
\item flexible, high-density electrode or nanosensor arrays;
\item low-noise analog front ends and energy-aware edge processing;
\item biocompatible packaging and reliable skin/interface contact.
\end{itemize}
These are requirements, not claims of a fabricated device. A valid hardware study must report material composition, sensor calibration, impedance, safety limits, sampling characteristics, and repeatability across users.

\section{Cognitive State Communication}
Rather than transmitting raw neural signals, NDAN transmits a state estimate and uncertainty:
\begin{equation}
 m_t = (\hat{z}_t, \Sigma_t, c_t, \pi_t),
\end{equation}
where $\hat{z}_t$ is the estimated state, $\Sigma_t$ its uncertainty, $c_t$ a consent scope, and $\pi_t$ a provenance record. An agent may act only when confidence and consent satisfy policy constraints. This design reduces unnecessary data exposure and makes uncertainty a first-class networking object.

\section{Evaluation Protocol}
The proposed validation sequence contains four stages:
\begin{enumerate}
\item synthetic signals for unit and integration testing;
\item public EEG/BCI datasets with subject-wise splits;
\item controlled laboratory sessions with informed consent;
\item hardware-in-the-loop evaluation of the sensing and edge pipeline.
\end{enumerate}
Metrics include balanced accuracy, calibration error, false activation rate, end-to-end latency, energy per inference, bandwidth per cognitive-state message, privacy leakage, and user-reported workload. No result should be generalized across users without subject-wise validation.

\begin{table}[h]
\centering
\caption{Required reporting matrix for future validation.}
\begin{tabular}{ll}
\toprule
Dimension & Required evidence \\
\midrule
Signal & Sampling, filtering, artifacts, sensor calibration \\
Model & Architecture, training split, uncertainty method \\
Network & Message size, latency, loss, adaptation policy \\
Safety & Consent, withdrawal, privacy, failure-safe behavior \\
Hardware & Materials, impedance, power, thermal limits \\
Reproducibility & Seeds, code, dataset license, environment \\
\bottomrule
\end{tabular}
\end{table}

\section{Governance and Neuro-Rights}
NDAN must not infer sensitive attributes without explicit purpose and consent. The governance layer applies data minimization, local processing where possible, cryptographic provenance, user-visible uncertainty, and an emergency stop. A cognitive estimate is not a diagnosis, legal fact, or authorization to act. Human-in-the-loop review is mandatory for consequential actions.

\section{Discussion}
NDAN connects nanoscale sensing, BCI signal processing, adaptive networks, and human-AI systems without collapsing them into an unsupported telepathy claim. The strongest near-term contribution is a reproducible systems architecture and evaluation framework. The most important research risk is domain shift: a model that works on synthetic signals may fail across users, electrode placements, motion artifacts, or changing tasks.

\section{Conclusion}
This paper defines a nanotechnology-ready architecture for NeuroDigital Interfaces and Adaptive Cognitive Networking. It provides a disciplined path from software prototypes to sensor and human-in-the-loop validation. Future work will implement public-dataset benchmarks, integrate the FCSTN/NDAN software lineage after its BCI test failures are corrected, and evaluate privacy-preserving cognitive-state communication.

\begin{thebibliography}{99}
\bibitem{wolpaw} J. R. Wolpaw et al., ``Brain-computer interfaces for communication and control,'' \emph{Clinical Neurophysiology}, 2002.
\bibitem{niedermeyer} E. Niedermeyer and F. Lopes da Silva, \emph{Electroencephalography}, 6th ed., 2011.
\bibitem{yuste} R. Yuste et al., ``Four ethical priorities for neurotechnologies and AI,'' \emph{Nature}, 2017.
\bibitem{ndandocs} F. M. Fuentes, ``NDAN and Adaptive Cognitive Networking,'' technical research package, 2026.
\end{thebibliography}
\end{document}
